I see three vulnerabilities in the headline result. First, the $p$-value's
framing. Under the declared
all-five rank-permutation frame, $p = 0.0167$ is the nominal rank-tail, exact
only under that combinatorial frame and not noise-aware. By the time the
five-behavior test set was fixed, two of the five effects
(\behavior{refusal}, \behavior{both\_sides}) were already known from the
replication gate. The ordering itself was committed before H2 separability,
overlap, or the three added target-cell outcomes were viewed, so those later
measurements did not determine the order. Two effects were therefore already
known when the final set was fixed, so this is not a fully fresh five-effect
$p$-value; no calibrated conditional $p$-value or selection model is
recorded. Second, seed noise versus ordering gaps: the two random-baseline
seeds differ by about 0.13 in \behavior{both\_sides} effect, while the
confirmed \behavior{validate\_feelings}-to-\behavior{bold\_formatting} gap is
0.047; every target cell is single-seed (pre-declared and budget-bound), so a
noise-aware ordering test might not reject where the exact permutation test
does. More target-training seeds are needed to test rank stability. Third,
what the index actually predicted: the measured separability values
mis-rank \behavior{validate\_feelings}, so the confirmed object is the
pre-registered ordering (written from priors and the post's figures
before the index values existed), and the index's own in-sample
monotonicity is partial. The defensible claim is ``the pre-registered
filterability ordering held under the nominal all-five rank-tail,'' not
``the index alone rank-orders filterability.''
